\documentclass[11pt]{article}

\usepackage[margin=1in]{geometry}
\usepackage{hyperref}
\usepackage{amsmath}

\title{CTC Beam Search and KenLM Reference}
\author{Danish ASR Project}
\date{\today}

\begin{document}

\maketitle

\section{Purpose}

This note summarizes the decoding references used for the OmniASR CTC
beam-search and KenLM methodology. The primary references are the ApX Machine
Learning notes on beam search with a language model and building n-gram models
with KenLM \cite{apxBeamSearchLm,apxKenlmNgram}. The Alexandra-proxy corpus is
based on Alexandra Institute's CoRal n-gram training code and public
ScandiWiki/ScandiReddit datasets
\cite{alexandraCoralNgram,alexandraScandiWiki,alexandraScandiReddit}.

\section{Beam Search}

Greedy CTC decoding chooses the most likely token at every acoustic timestep,
then applies CTC cleanup by collapsing repeated tokens and removing blank
tokens. It is fast, but it commits to local choices immediately.

Beam search keeps multiple candidate transcripts alive. At each timestep, the
decoder extends the current hypotheses, scores the extensions, sorts them, and
keeps only the top candidates. The number of retained candidates is the beam
width. A beam width of 64 means that up to 64 partial transcripts survive each
pruning step.

Without a language model, beam search is still an acoustic-only search. It can
recover from local ambiguity better than greedy decoding, but it does not know
which word sequences are more natural.

\section{Language-Model Scoring}

When a language model is added, the decoder combines acoustic and language
evidence. A common scoring form is:

\[
\mathrm{score}(W) =
\mathrm{score}_{\mathrm{acoustic}}(W)
+ \alpha \, \mathrm{score}_{\mathrm{LM}}(W)
+ \beta \, \mathrm{word\_count}(W)
\]

Here, $\alpha$ controls the language-model weight and $\beta$ is a word
insertion bonus. The word insertion term helps avoid an excessive preference
for very short hypotheses.

\section{KenLM}

KenLM is not a variant of beam search. Beam search is the search algorithm;
KenLM is a toolkit for building and querying n-gram language models. An n-gram
language model estimates the probability of a word using only the previous
$n-1$ words. For a trigram model:

\[
P(w_t \mid w_{t-2}, w_{t-1})
\]

During beam search, the decoder can query KenLM whenever a candidate hypothesis
completes a word. The KenLM score then influences whether that hypothesis
survives pruning.

\section{Project Usage}

In this project:

\begin{itemize}
  \item \texttt{CTC no\_lm} means greedy CTC decoding.
  \item \texttt{beam} means CTC beam search without an external language model.
  \item \texttt{CTC LM-enabled} means CTC beam search with a KenLM n-gram model.
\end{itemize}

The main Alexandra-proxy KenLM artifact is a 3-gram model trained on Danish
ScandiWiki and Danish ScandiReddit text, with CoRal-v3 test transcripts
excluded to avoid evaluation leakage. A smaller CoRal-v3-train-only LM remains
available as an ablation, but should not be described as Alexandra-aligned. The
current documented proxy settings are:

\begin{itemize}
  \item beam width: 64
  \item alpha: 0.5
  \item beta: 1.5
\end{itemize}

If alpha and beta are tuned, they should be tuned only on a validation split and
then frozen before evaluating on the CoRal-v3 test split.

\bibliographystyle{plain}
\bibliography{ctc-beam-kenlm-reference}

\end{document}
